\begin{theorem}[о локальной структуре отображения]\label{local-structure}
	Пусть выполнены следующие условия:
	\begin{enumerate}
		\item $G \subset \Cm$ "--- область
		\item Функция $f$ регулярна на $G$
		\item $z_0 \in G$, $w_0 := f(z_0)$
		\item Для некоторого $n \in \mathbb{N} \; : n \ge 2$ выполнено, что $f'(z_0) = \ldots = f^{(n - 1)}(z_0) = 0$, но $f^{(n)}(z_0) \ne 0$
	\end{enumerate}
	
	Тогда $\exists \; \delta > 0$ и $\epsilon > 0$ такие, что для $ \forall \; \widetilde w \in \mathring B_\epsilon(w_0)$ уравнение $f(z) = \widetilde w$ имеет ровно $n$ различных решений на $B_\delta(z_0)$.
\end{theorem}

\begin{proof}
	Пусть сначала $w_0 = 0$. Тогда точка $z_0$ "--- изолированный ноль функций $f$ и $f'$, поскольку $f^{(n)}(z_0) \ne 0$ (если предположить противное, получим, что $f \equiv 0$ по теореме о единственности). Выберем $\delta > 0$ такое, что других нулей на $\overline{B_\delta(z_0)}$ у этих функций нет. Положим $\epsilon = \min\{|f(z)| : |z - z_0| = \delta\}$. Зафиксируем $\widetilde w \in \mathring B_\epsilon(w_0)$ и проверим условия теоремы Руше для функций $f$ и $-\widetilde{w}$ на области $B_\delta(z_0)$:
	\[|f(z)| \ge \epsilon > |\!-\!\widetilde{w}|,~|z - z_0| = \delta\]
	
	Значит, уравнение $f(z) = \widetilde w$, как и $f(z) = 0$, имеет $n$ корней на $B_\delta(z_0)$ с учетом кратности. Очевидно, что $f$ имеет ровно 1 корень кратности $n$,  так как первые $n - 1$ производные $= 0$. \\Осталось доказать, что все нули функции $f(z) - \widetilde w$ имеют порядок $1$, это правда поскольку \\ $(f- \widetilde w)' \ne 0$ на $\mathring B_\delta(z_0)$ (в силу того, что $z_0$ - изолированный ноль $f'$). Если же $w_0 \ne 0$, то доказанную часть теоремы можно применить к функции $f - w_0$.
\end{proof}

\begin{note}
	Таким образом, на $\mathring B_\epsilon(w_0)$ определены функции $g_1, \dotsc, g_n$, не совпадающие ни в одной точке, такие, что для каждого $k \in \{1, \dotsc, n\}$ на $\mathring B_\epsilon(w_0)$ выполнено тождество $f(g_k(w)) \equiv w$. Кроме того, вместо пары $\delta$ и $\epsilon$ в доказательстве можно выбрать любые пары $\delta' \in (0, \delta)$ и $\epsilon' := \min\{|f(z)| : |z - z_0| = \delta'\}$, поэтому для каждого $k \in \{1, \dotsc, n\}$ выполнено $\lim\limits_{w \to w_0}g_k(w) \equiv z_0$.
\end{note}

\begin{theorem}[принцип сохранения области]
	Пусть непостоянная функция $f$ регулярна на области $G$. Тогда $f(G)$ "--- тоже область.
\end{theorem}

\begin{proof}
	Докажем, что множество $f(G)$ "--- открытое. Пусть $w_0 \in f(G)$ тогда \\ $\exists \; z_0 \in G \; : \; f(z_0) = w_0$. Если $f'(z_0) \ne 0$, то по теореме об обратной функции $\exists \; \delta > 0$ и $\epsilon > 0$ такие, что  $\forall \; w \in B_\epsilon(w_0)$ уравнение $f(z) = w$ имеет единственное решение на $B_\delta(z_0)$, поэтому $B_\epsilon(w_0) \subset f(G)$. Если же, наоборот, $f'(z_0) = 0$, то, поскольку $f$ непостоянна, $ \exists \; n \in \mathbb{N} \; : n \ge 2$ такое, что $f^{(n)}(z_0) \ne 0$, и применима теорема о локальной структуре отображения, дающая аналогичный результат.
	
	Теперь докажем связность множества $f(G)$. Пусть $w_1, w_2 \in f(G)$ тогда для некоторых $z_1, z_2 \in G$ выполнено $f(z_1) = w_1$ и $f(z_2) = w_2$. Пусть $\gamma$ "--- кривая, соединяющая точки $z_1$ и $z_2$, такая что $M(\gamma) \subset G$. Тогда кривая $\Gamma = f(\gamma)$ соединяет точки $w_1$ и $w_2$, причем $M(\Gamma) \subset f(G)$. Получено требуемое.
\end{proof}
